
\documentclass{article}
\usepackage{colortbl}
\usepackage{makecell}
\usepackage{multirow}
\usepackage{supertabular}

\begin{document}

\newcounter{utterance}

\centering \large Interaction Transcript for game `ifeval', experiment `ifeval', episode 429 with refineonly{-}Qwen3.5{-}27B.
\vspace{24pt}

{ \footnotesize  \setcounter{utterance}{1}
\setlength{\tabcolsep}{0pt}
\begin{supertabular}{c@{$\;$}|p{.15\linewidth}@{}p{.15\linewidth}p{.15\linewidth}p{.15\linewidth}p{.15\linewidth}p{.15\linewidth}}
    \# & \multicolumn{2}{c}{Player} && \multicolumn{2}{c}{Game Master} \\
    \hline

    \theutterance \stepcounter{utterance}  
    & & & \multicolumn{4}{p{0.6\linewidth}}{
        \cellcolor[rgb]{0.9,0.9,0.9}{
            \makecell[{{p{\linewidth}}}]{
                \texttt{\tiny{[P1$\langle$GM]}}
                \texttt{Is "jiang" a Chinese name? What are other names similar to "jiang"? Separate your two answers with ******} \\
            }
        }
    }
     \\ \\

    \theutterance \stepcounter{utterance}  
    & \multicolumn{4}{p{0.6\linewidth}}{
        \cellcolor[rgb]{0.561,0.737,0.561}{
            \makecell[{{p{\linewidth}}}]{
                \texttt{\tiny{[P1$\rangle$GM]}}
                \texttt{Yes, "Jiang" is a very common Chinese surname (family name). It is written as **江** in Chinese characters, which literally means "river." It is one of the most frequent surnames in China, particularly prevalent in the southern regions, and is also found among Chinese communities worldwide. While it is primarily a surname, it can occasionally appear as a given name, though this is less common than its use as a family name.} \\
\\ 
\texttt{******} \\
\\ 
\texttt{Other Chinese names similar to "Jiang" can be categorized by their pronunciation, meaning, or character structure:} \\
\\ 
\texttt{*   **Similar Pronunciation (Pinyin):**} \\
\texttt{    *   **Jiang (姜):** A different surname written with a different character (meaning "ginger"), but pronounced exactly the same.} \\
\texttt{    *   **Jiang (将):** Another surname meaning "general" or "to lead," also pronounced the same.} \\
\texttt{    *   **Qiang (强):** Meaning "strong," which shares the "ang" ending and a similar initial consonant sound.} \\
\texttt{    *   **Jian (建/健):** Meaning "build" or "healthy," sharing the "J" sound and similar vowel structure.} \\
\\ 
\texttt{*   **Similar Meaning (Water/Nature):**} \\
\texttt{    *   **He (河):** Meaning "river," often used as a surname or given name.} \\
\texttt{    *   **Hao (浩):** Meaning "vast" or "grand," often used in given names and containing the water radical (氵) like 江.} \\
\texttt{    *   **Tao (涛):** Meaning "waves," also containing the water radical.} \\
\texttt{    *   **Yun (云):** Meaning "cloud," often paired with water{-}related names in poetry and naming conventions.} \\
\\ 
\texttt{*   **Common Surnames of Similar Frequency:**} \\
\texttt{    *   **Chen (陈):** One of the most common surnames in China, often listed alongside Jiang in frequency rankings.} \\
\texttt{    *   **Liu (刘):** A very common surname, frequently cited in similar contexts of prevalence.} \\
\texttt{    *   **Wang (王):** The most common surname in China, often grouped with Jiang as a top{-}tier family name.} \\
            }
        }
    }
    & & \\ \\

    \theutterance \stepcounter{utterance}  
    & & & \multicolumn{2}{p{0.3\linewidth}}{
        \cellcolor[rgb]{0.9,0.9,0.9}{
            \makecell[{{p{\linewidth}}}]{
                \texttt{\tiny{[GM$|$GM]}}
                {'combination:two_responses': {}}
            }
        }
    }
    & & \\ \\

    \theutterance \stepcounter{utterance}  
    & & & \multicolumn{2}{p{0.3\linewidth}}{
        \cellcolor[rgb]{0.9,0.9,0.9}{
            \makecell[{{p{\linewidth}}}]{
                \texttt{\tiny{[GM$|$GM]}}
                \texttt{game\_result = WIN} \\
            }
        }
    }
    & & \\ \\

\end{supertabular}
}

\end{document}
